\documentclass{article}
\usepackage[margin=0.8in]{geometry}
\usepackage{amsmath,amssymb,amsthm,mathtools}
\usepackage{mathrsfs,bm,bbm}
\usepackage{graphicx,booktabs,array,multirow,tabularx,longtable}
\usepackage{enumitem}
\usepackage{xcolor}
\usepackage{algorithm,algpseudocode}
\usepackage{subcaption,tikz}
\usepackage{siunitx,cancel,accents}
\usepackage{microtype}
\usepackage[numbers,sort&compress]{natbib}
\usepackage[colorlinks=true,linkcolor=blue,citecolor=blue,urlcolor=blue]{hyperref}
\usepackage{cleveref}
\setlength{\emergencystretch}{2em}
\newtheorem{assumption}{Assumption}
\newtheorem{definition}{Definition}
\newtheorem{theorem}{Theorem}
\newtheorem{lemma}{Lemma}
\newtheorem{proposition}{Proposition}
\newtheorem{openendedquestion}{Open-ended Question}
\newtheorem{conjecture}{Conjecture}
\newcommand{\coreref}[1]{\ref{#1}}

\begin{document}

\sloppy

\section*{scm\_\allowbreak{}supportface\_\allowbreak{}complete\_\allowbreak{}id}

\noindent\textit{Specialization:} v1\par

\paragraph{TL;DR.} Exact structural zeros are identifying information inside the canonical quasi-Markovian class \(\Theta_S\), which has one mutually independent finite response root per bidirected-clique district. First-forbidden-prefix saturation gives its unique maximal Cartesian response support, and positive-prefix telescoping recovers linear district signatures. On \(\Theta_S\), a deterministic intervention query is uniformly identified exactly when its response indicator belongs to the tensor product of those signature row spaces: membership yields one denominator-safe rational functional, while nonmembership yields strictly positive rational same-support, same-law twins with different causal answers. Because those twins remain available in every semantics-preserving ambient class \(\mathcal M_S\supseteq\Theta_S\), identification in such a class implies tensor membership; the success direction is not transferred beyond \(\Theta_S\). The rational functional has fixed-face multinomial influence-function inference, while complete support-scoped do-search compilation remains open.

\section{Project justification}

\paragraph{Gap.} Existing structural-zero work gives positivity checks or selected supported operations, quasi-Markovian credal work gives compatible-set representations and numerical bounds, and broader canonical discrete-SCM work gives polynomial optimization. None of these supplies a support-only formula-or-exact-interior-twins criterion that is necessary and sufficient uniformly over the canonical quasi-Markovian face \(\Theta_S\).

\paragraph{Niche.} Within \(\Theta_S\), independent district roots turn observational zeros into a maximal product response support and supported prefix ratios into separately identified linear signatures. This yields finite tensor geometry. Its failure certificate transfers to every semantics-preserving supermodel class \(\mathcal M_S\supseteq\Theta_S\), whereas its rational success certificate remains class-specific.

\paragraph{Fill.} The project proves support saturation, district-signature recovery, the tensor iff on \(\Theta_S\), a denominator-safe rational formula, strictly positive rational same-law twins, and fixed-face multinomial inference. It also records the asymmetric ambient-class consequence: identification over \(\mathcal M_S\supseteq\Theta_S\) requires tensor membership, but no success theorem is claimed for Duarte's broader overlapping-latent class.

\section{Related work}

Hwang, Choe, Kwon, and Lee motivate supported \(Q\)-factor operations when strict positivity fails, while complete data-adaptive identification remains unresolved there. Zaffalon, Antonucci, and Cabañas and their 2024 extension provide the closest class-level substrate: quasi-Markovian observational constraints become separate rootwise credal constraints and compatible models admit credal-network analysis. This paper's delta within that canonical quasi-Markovian setting \(\Theta_S\) is the maximal saturated product support and the uniform symbolic formula-or-strictly-positive-rational-twins dichotomy. Duarte, Finkelstein, Knox, Mummolo, and Shpitser justify finite canonical response types and exact polynomial optimization in a broader discrete-SCM framework that may include overlapping latent-parent neighborhoods. The present success theorem is not extended to that broader class. Only the converse transfers: the twins already lie in \(\Theta_S\), so they also refute identification in every semantics-preserving ambient class containing it. Shridharan and Iyengar motivate districtwise multilinear reduction; Chen and Darwiche motivate constrained identification by invariants and distinguish cardinality information from exact support; Tikka, Hyttinen, and Karvanen provide the search-rule framework for the open compiler; and Bhattacharya, Nabi, and Shpitser provide the differentiation template for inference from an identified graphical functional.

\section{Setup}

Target (estimand / structural parameter): \(\tau_{x,y}(\theta)=P_{\theta}(Y=y\mid\operatorname{do}(X=x))\).
Primary functional / estimator / diagnostic: \(F_{\lambda}(P)=\sum_{a\in(\{0\}\cup S)^m}\lambda[a]\prod_{j=1}^m t_j(P)[a_j]\), where \(\lambda B=h_{x,y}|_{T(S)}\).
\begin{itemize}
  \item \texttt{n} --- \texttt{positive integer}
  \par\smallskip\noindent\emph{Definition.} \texttt{The number of observed vertices.}
  \item \texttt{V} --- \texttt{ordered finite vertex set}; \(V=\{V_1,\ldots,V_n\}\); \texttt{Observed variables.}
  \item \texttt{prec} --- \texttt{topological order}
  \par\smallskip\noindent\emph{Definition.} \(V_1\prec\cdots\prec V_n\).
  \item \texttt{G} --- \texttt{acyclic directed mixed graph}; A graph on \(V\) whose bidirected districts are cliques.; \texttt{Known causal graph.}
  \item \texttt{m} --- \texttt{positive integer}
  \par\smallskip\noindent\emph{Definition.} The number of bidirected districts of \(G\).
  \item \texttt{D} --- partition of \(V\); \(D=(D_1,\ldots,D_m)\)
  \par\smallskip\noindent\emph{Definition.} The bidirected districts of \(G\).
  \item \texttt{X\_\allowbreak{}i} --- \texttt{finite set}; \(\mathcal X_i\)
  \par\smallskip\noindent\emph{Definition.} The state space of \(V_i\).
  \item \texttt{X\_\allowbreak{}V} --- \texttt{finite product set}; \(\mathcal X_V=\prod_{i=1}^n\mathcal X_i\)
  \par\smallskip\noindent\emph{Definition.} \texttt{The observational cell space.}
  \item \texttt{S} --- \texttt{nonempty subset}; \(S\subseteq\mathcal X_V\); \texttt{Supplied realizable exact observational support.}
  \item \texttt{S\_\allowbreak{}le\_\allowbreak{}k} --- \texttt{prefix-\allowbreak{}support family}; \(S^{\le k}\subseteq\prod_{i=1}^k\mathcal X_i\)
  \par\smallskip\noindent\emph{Definition.} \(S^{\le k}=\{s_{\le k}:s\in S\}\), with \(S^{\le0}=\{()\}\).
  \item \texttt{R\_\allowbreak{}j} --- \texttt{finite district response space}; \(\mathcal R_j=\prod_{i\in D_j}\mathcal X_i^{\mathcal X_{\operatorname{pa}_G(i)}}\)
  \par\smallskip\noindent\emph{Definition.} A member is a tuple of deterministic response tables for all vertices in district \(D_j\).
  \item \texttt{theta} --- \texttt{canonical response-\allowbreak{}law parameter}; \(\theta=(\theta_1,\ldots,\theta_m)\in\prod_{j=1}^m\Delta(\mathcal R_j)\)
  \par\smallskip\noindent\emph{Definition.} The joint root law is the product \(\bigotimes_{j=1}^m\theta_j\).
  \item \texttt{theta\_\allowbreak{}j} --- \texttt{district probability mass function}; \(\theta_j\in\Delta(\mathcal R_j)\)
  \par\smallskip\noindent\emph{Definition.} The \(j\)-th component of \(\theta\).
  \item \texttt{r\_\allowbreak{}j} --- \texttt{district response type}; \(r_j\in\mathcal R_j\)
  \item \texttt{r} --- \texttt{joint response tuple}; \(r=(r_1,\ldots,r_m)\in\prod_{j=1}^m\mathcal R_j\)
  \item \texttt{X} --- \texttt{vertex subset}; \(X\subseteq V\); \texttt{Intervened variables.}
  \item \texttt{x} --- \texttt{intervention value}; \(x\in\mathcal X_X\)
  \item \texttt{Y} --- \texttt{vertex subset}; \(Y\subseteq V\setminus X\); \texttt{Outcome variables.}
  \item \texttt{y} --- \texttt{outcome value}; \(y\in\mathcal X_Y\)
  \item \texttt{ev\_\allowbreak{}x} --- \texttt{deterministic evaluation map}; \(\operatorname{ev}_x:\prod_{j=1}^m\mathcal R_j\to\mathcal X_V\)
  \par\smallskip\noindent\emph{Definition.} Evaluate the response tables in order \(\prec\), replacing the equations for \(X\) by the constants \(x\); \(\operatorname{ev}_{\varnothing}\) denotes factual evaluation.
  \item \texttt{P\_\allowbreak{}theta} --- \texttt{observational probability mass function}; \(P_{\theta}\in\Delta(\mathcal X_V)\)
  \par\smallskip\noindent\emph{Definition.} \(P_{\theta}(v)=\sum_{r\in\prod_j\mathcal R_j}\prod_{j=1}^m\theta_j(r_j)\mathbf 1\{\operatorname{ev}_{\varnothing}(r)=v\}\).
  \item \texttt{Theta\_\allowbreak{}S} --- \texttt{model class}; \(\Theta_S\subseteq\prod_{j=1}^m\Delta(\mathcal R_j)\)
  \par\smallskip\noindent\emph{Definition.} Defined by \(\mathrm{def:model-class}\).
  \item \texttt{h\_\allowbreak{}xy} --- \texttt{binary response tensor}; \(h_{x,y}\in\{0,1\}^{\prod_j\mathcal R_j}\)
  \par\smallskip\noindent\emph{Definition.} \(h_{x,y}(r)=\mathbf 1\{\operatorname{ev}_x(r)_Y=y\}\).
  \item \texttt{tau\_\allowbreak{}theta} --- \texttt{causal probability}; \([0,1]\)
  \par\smallskip\noindent\emph{Definition.} \(\tau_{x,y}(\theta)=\sum_{r\in\prod_j\mathcal R_j}h_{x,y}(r)\prod_{j=1}^m\theta_j(r_j)\).
  \item \texttt{F\_\allowbreak{}i} --- \texttt{forbidden-\allowbreak{}prefix set}; \(\mathcal F_i(S)\)
  \par\smallskip\noindent\emph{Definition.} Defined by \(\mathrm{def:forbidden-prefixes}\).
  \item \texttt{b\_\allowbreak{}j\_\allowbreak{}k} --- \texttt{binary local prefix indicator}; \(b_j^k(w,r_j)\in\{0,1\}\)
  \par\smallskip\noindent\emph{Definition.} Defined by \(\mathrm{def:prefix-indicator}\).
  \item \texttt{T\_\allowbreak{}j} --- \texttt{support-\allowbreak{}restricted response space}; \(T_j(S)\subseteq\mathcal R_j\)
  \par\smallskip\noindent\emph{Definition.} Defined by \(\mathrm{def:saturated-local-support}\).
  \item \texttt{T} --- \texttt{Cartesian response support}; \(T(S)=\prod_{j=1}^mT_j(S)\)
  \par\smallskip\noindent\emph{Definition.} Defined by \(\mathrm{def:saturated-product-support}\).
  \item \texttt{P} --- \texttt{observational probability mass function}; \(P\in\operatorname{relint}(\Delta(S))\), where \(\Delta(S)=\{p\in\Delta(\mathcal X_V):p(v)=0\ \text{for }v\notin S\}\); \texttt{Generic law on the relative interior of the supplied face.}
  \item \texttt{B\_\allowbreak{}j} --- \texttt{rational matrix}; \(B_j\in\{0,1\}^{(\{0\}\cup S)\times T_j(S)}\)
  \par\smallskip\noindent\emph{Definition.} Defined by \(\mathrm{def:district-signature-matrix}\).
  \item \texttt{t\_\allowbreak{}j} --- \texttt{observable district-\allowbreak{}signature vector}; \(t_j(P)\in\mathbb R^{\{0\}\cup S}\)
  \par\smallskip\noindent\emph{Definition.} Defined by \(\mathrm{def:observable-signature}\).
  \item \texttt{B} --- \texttt{tensor signature matrix}; \(B=\bigotimes_{j=1}^mB_j\)
  \par\smallskip\noindent\emph{Definition.} Defined by \(\mathrm{def:tensor-signature-matrix}\).
  \item \texttt{lambda} --- \texttt{rational row coefficient vector}; \(\lambda\in\mathbb Q^{(\{0\}\cup S)^m}\); A solution of \(\lambda B=h_{x,y}|_{T(S)}\), when one exists.
  \item \texttt{F\_\allowbreak{}lambda} --- \texttt{rational observational functional}; \(F_{\lambda}:\operatorname{relint}(\Delta(S))\to\mathbb R\)
  \par\smallskip\noindent\emph{Definition.} Defined by \(\mathrm{def:rational-functional}\).
  \item \texttt{D\_\allowbreak{}S} --- \texttt{common denominator polynomial}; \(D_S:\Delta(S)\to\mathbb R\)
  \par\smallskip\noindent\emph{Definition.} Defined by \(\mathrm{def:common-denominator}\).
  \item \texttt{epsilon} --- \texttt{positive constant}; \(\epsilon>0\); \texttt{Fixed distance from the boundary of the support face.}
  \item \texttt{sigma\_\allowbreak{}min} --- \texttt{positive constant}; \(\underline\sigma>0\); \texttt{Uniform lower bound for standard deviations.}
  \item \texttt{N} --- \texttt{positive integer}
  \par\smallskip\noindent\emph{Definition.} \texttt{The observational sample size.}
  \item \texttt{O\_\allowbreak{}k} --- \texttt{observed categorical draw}; \(O_k\in S\)
  \par\smallskip\noindent\emph{Definition.} The \(k\)-th observational draw, for \(k\in\{1,\ldots,N\}\).
  \item \texttt{C\_\allowbreak{}N} --- \texttt{multinomial count vector}; \(C_N=(C_N(s):s\in S)\)
  \par\smallskip\noindent\emph{Definition.} \(C_N(s)=\sum_{k=1}^N\mathbf 1\{O_k=s\}\).
  \item \texttt{P\_\allowbreak{}hat\_\allowbreak{}N} --- \texttt{empirical probability mass function}; \(\widehat P_N\in\Delta(S)\)
  \par\smallskip\noindent\emph{Definition.} \(\widehat P_N(s)=C_N(s)/N\).
  \item \texttt{tau\_\allowbreak{}hat\_\allowbreak{}N} --- \texttt{plug-\allowbreak{}in estimator}; \(\widehat\tau_N\in\mathbb R\)
  \par\smallskip\noindent\emph{Definition.} Defined by \(\mathrm{def:plugin-estimator}\).
  \item \texttt{phi\_\allowbreak{}P} --- \texttt{face-\allowbreak{}tangent influence function}; \(\phi_P:S\to\mathbb R\)
  \par\smallskip\noindent\emph{Definition.} Defined by \(\mathrm{def:influence-function}\).
  \item \texttt{sigma2\_\allowbreak{}P} --- \texttt{asymptotic variance}; \(\sigma^2(P)\in[0,\infty)\)
  \par\smallskip\noindent\emph{Definition.} \(\sigma^2(P)=\sum_{s\in S}P(s)\phi_P(s)^2\).
  \item \texttt{P\_\allowbreak{}face} --- \texttt{fixed interior statistical class}; \(\mathcal P_{S,\epsilon,\underline\sigma}\)
  \par\smallskip\noindent\emph{Definition.} Defined by \(\mathrm{def:fixed-interior-class}\).
  \item \texttt{Z} --- \texttt{binary observed variable}; \(\{0,1\}\); \texttt{Instrument in the five-\allowbreak{}cell control.}
  \item \texttt{G\_\allowbreak{}bow} --- \texttt{fixed acyclic directed mixed graph}
  \par\smallskip\noindent\emph{Definition.} The binary bow graph \(X\to Y\) and \(X\leftrightarrow Y\).
  \item \texttt{S\_\allowbreak{}diag} --- \texttt{fixed support}
  \par\smallskip\noindent\emph{Definition.} \(S_{\mathrm{diag}}=\{(0,0),(1,1)\}\) for \((X,Y)\).
  \item \texttt{G\_\allowbreak{}IV} --- \texttt{fixed acyclic directed mixed graph}
  \par\smallskip\noindent\emph{Definition.} The binary instrumental-variable graph \(Z\to X\to Y\) and \(X\leftrightarrow Y\).
  \item \texttt{S\_\allowbreak{}IV5} --- \texttt{fixed support}
  \par\smallskip\noindent\emph{Definition.} \(S_{\mathrm{IV5}}=\{(0,0,0),(0,1,0),(1,0,0),(1,0,1),(1,1,1)\}\) for \((Z,X,Y)\).
  \item \texttt{M\_\allowbreak{}S} --- \texttt{ambient finite-\allowbreak{}SCM class}; \(\mathcal M_S\)
  \par\smallskip\noindent\emph{Definition.} A class of finite acyclic SCMs with factual support \(S\) on the same observed graph and state spaces as \(\Theta_S\), equipped with observational and intervention semantics whose restrictions to embedded canonical models are \(P_\theta\) and \(\tau_{x,y}(\theta)\).
\end{itemize}

\section{Assumptions}

\begin{assumption}[ass:exact-support]\label{ass:exact-support}
\(\operatorname{supp}(P_{\theta})=S\).
\par\smallskip\noindent\emph{Standard} (Exact observational compatibility constraint, \cite{ChenDarwiche2024Constrained}).
\end{assumption}

\begin{assumption}[ass:realizable-support]\label{ass:realizable-support}
\(\Theta_S\ne\varnothing\).
\par\smallskip\noindent\emph{Standard} (Feasibility of the canonical discrete SCM constraints, \cite{DuarteFinkelsteinKnoxMummoloShpitser2023}).
\end{assumption}

\begin{assumption}[ass:iid-multinomial-sampling]\label{ass:iid-multinomial-sampling}
\(C_N\sim\operatorname{Multinomial}(N,(P(s))_{s\in S})\).
\par\smallskip\noindent\emph{Standard} (Independent multinomial sampling, \cite{vanDerVaart1998}).
\end{assumption}

\begin{assumption}[ass:fixed-face-interior]\label{ass:fixed-face-interior}
\(\min_{s\in S}P(s)\ge\epsilon\).
\par\smallskip\noindent\emph{Standard} (Fixed multinomial interior condition, \cite{vanDerVaart1998}).
\end{assumption}

\begin{assumption}[ass:nondegenerate-variance]\label{ass:nondegenerate-variance}
\(\sigma^2(P)\ge\underline\sigma^2\).
\par\smallskip\noindent\emph{Standard} (Uniform variance nondegeneracy, \cite{vanDerVaart1998}).
\end{assumption}

\section{Definitions}

\begin{definition}[def:model-class]\label{def:model-class}
\par\smallskip\noindent\emph{Defined object.} \texttt{Canonical quasi-\allowbreak{}Markovian exact-\allowbreak{}support SCM class}
\par\smallskip\noindent\emph{Construction.} The canonical quasi-Markovian exact-support class is
\[
\Theta_S=\left\{\theta\in\prod_{j=1}^m\Delta(\mathcal R_j):
\operatorname{supp}(P_\theta)=S\right\}.
\]
Its product parameterization is the one-independent-finite-response-root-per-bidirected-clique-district parameterization fixed by \(\theta\).
\par\smallskip\noindent\emph{Member properties.} \texttt{ass:\allowbreak{}exact-\allowbreak{}support}.
\end{definition}

\begin{definition}[def:forbidden-prefixes]\label{def:forbidden-prefixes}
\par\smallskip\noindent\emph{Defined object.} \texttt{First forbidden extensions}
\par\smallskip\noindent\emph{Construction.} \(\mathcal F_i(S)=\{(w,a):w\in S^{\le i-1},\ a\in\mathcal X_i,\ wa\notin S^{\le i}\}\) for every \(i\in\{1,\ldots,n\}\).
\par\smallskip\noindent\emph{Inputs.} \texttt{S}; \texttt{prec}.
\end{definition}

\begin{definition}[def:prefix-indicator]\label{def:prefix-indicator}
\par\smallskip\noindent\emph{Defined object.} \texttt{District prefix indicator}
\par\smallskip\noindent\emph{Construction.} \(b_j^k(w,r_j)=\prod_{\ell\le k:\,V_\ell\in D_j}\mathbf 1\{r_{j,\ell}(w_{\operatorname{pa}_G(V_\ell)})=w_\ell\}\) for \(j\in\{1,\ldots,m\}\), \(k\in\{0,\ldots,n\}\), \(w\in\prod_{\ell=1}^k\mathcal X_\ell\), and \(r_j\in\mathcal R_j\), with the empty product equal to one.
\par\smallskip\noindent\emph{Inputs.} \texttt{G}; \texttt{prec}; \texttt{D}; \texttt{r\_\allowbreak{}j}.
\end{definition}

\begin{definition}[def:saturated-local-support]\label{def:saturated-local-support}
\par\smallskip\noindent\emph{Defined object.} \texttt{First-\allowbreak{}forbidden-\allowbreak{}prefix saturated local support}
\par\smallskip\noindent\emph{Construction.} \(T_j(S)=\{r_j\in\mathcal R_j:b_j^i(wa,r_j)=0\ \text{for every }i\in\{1,\ldots,n\}\text{ with }V_i\in D_j\text{ and every }(w,a)\in\mathcal F_i(S)\}\) for every \(j\in\{1,\ldots,m\}\).
\par\smallskip\noindent\emph{Inputs.} \texttt{S}; \texttt{G}; \texttt{prec}; \texttt{D}.
\end{definition}

\begin{definition}[def:saturated-product-support]\label{def:saturated-product-support}
\par\smallskip\noindent\emph{Defined object.} \texttt{Saturated Cartesian response support}
\par\smallskip\noindent\emph{Construction.} \(T(S)=\prod_{j=1}^mT_j(S)\).
\par\smallskip\noindent\emph{Inputs.} \texttt{T\_\allowbreak{}j}; \texttt{m}.
\end{definition}

\begin{definition}[def:district-signature-matrix]\label{def:district-signature-matrix}
\par\smallskip\noindent\emph{Defined object.} \texttt{Supported district signature matrix}
\par\smallskip\noindent\emph{Construction.} For rows \(a\in\{0\}\cup S\) and columns \(r_j\in T_j(S)\), set \(B_j[0,r_j]=1\) and \(B_j[s,r_j]=b_j^n(s,r_j)\) for \(s\in S\).
\par\smallskip\noindent\emph{Inputs.} \texttt{S}; \texttt{T\_\allowbreak{}j}; \texttt{b\_\allowbreak{}j\_\allowbreak{}k}.
\end{definition}

\begin{definition}[def:observable-signature]\label{def:observable-signature}
\par\smallskip\noindent\emph{Defined object.} \texttt{Positive-\allowbreak{}prefix observable district signature}
\par\smallskip\noindent\emph{Construction.} Set \(t_j(P)[0]=1\) and \(t_j(P)[s]=\prod_{i:\,V_i\in D_j}P(V_{\le i}=s_{\le i})/P(V_{<i}=s_{<i})\) for every \(s\in S\) and \(j\in\{1,\ldots,m\}\).
\par\smallskip\noindent\emph{Inputs.} \texttt{P}; \texttt{S}; \texttt{D}; \texttt{prec}.
\end{definition}

\begin{definition}[def:tensor-signature-matrix]\label{def:tensor-signature-matrix}
\par\smallskip\noindent\emph{Defined object.} \texttt{Tensor district signature matrix}
\par\smallskip\noindent\emph{Construction.} \(B=\bigotimes_{j=1}^mB_j\), with columns indexed by \(T(S)\) and rows indexed by \((\{0\}\cup S)^m\).
\par\smallskip\noindent\emph{Inputs.} \texttt{B\_\allowbreak{}j}; \texttt{m}.
\end{definition}

\begin{definition}[def:rational-functional]\label{def:rational-functional}
\par\smallskip\noindent\emph{Defined object.} \texttt{Compiled rational functional}
\par\smallskip\noindent\emph{Construction.} \(F_{\lambda}(P)=\sum_{(a_1,\ldots,a_m)\in(\{0\}\cup S)^m}\lambda[a_1,\ldots,a_m]\prod_{j=1}^m t_j(P)[a_j]\).
\par\smallskip\noindent\emph{Inputs.} \texttt{lambda}; \texttt{t\_\allowbreak{}j}; \texttt{P}.
\end{definition}

\begin{definition}[def:common-denominator]\label{def:common-denominator}
\par\smallskip\noindent\emph{Defined object.} \texttt{Face-\allowbreak{}positive common denominator}
\par\smallskip\noindent\emph{Construction.} \(D_S(P)=\prod_{i=1}^n\prod_{w\in S^{\le i-1}}P(V_{<i}=w)\), where the empty-prefix factor equals one.
\par\smallskip\noindent\emph{Inputs.} \texttt{P}; \texttt{S}; \texttt{prec}.
\end{definition}

\begin{definition}[def:plugin-estimator]\label{def:plugin-estimator}
\par\smallskip\noindent\emph{Defined object.} \texttt{Support-\allowbreak{}face plug-\allowbreak{}in estimator}
\par\smallskip\noindent\emph{Construction.} \(\widehat\tau_N=F_{\lambda}(\widehat P_N)\) when \(D_S(\widehat P_N)>0\), and \(\widehat\tau_N=0\) otherwise.
\par\smallskip\noindent\emph{Inputs.} \texttt{lambda}; \texttt{P\_\allowbreak{}hat\_\allowbreak{}N}.
\end{definition}

\begin{definition}[def:influence-function]\label{def:influence-function}
\par\smallskip\noindent\emph{Defined object.} \texttt{Saturated-\allowbreak{}face influence function}
\par\smallskip\noindent\emph{Construction.} \(\phi_P(s)=\partial F_{\lambda}(P)/\partial P(s)-\sum_{u\in S}P(u)\,\partial F_{\lambda}(P)/\partial P(u)\) for \(s\in S\).
\par\smallskip\noindent\emph{Inputs.} \texttt{F\_\allowbreak{}lambda}; \texttt{P}.
\end{definition}

\begin{definition}[def:fixed-interior-class]\label{def:fixed-interior-class}
\par\smallskip\noindent\emph{Defined object.} \texttt{Fixed epsilon-\allowbreak{}interior multinomial class}
\par\smallskip\noindent\emph{Construction.} \(\mathcal P_{S,\epsilon,\underline\sigma}=\{P\in\operatorname{relint}(\Delta(S)):\min_{s\in S}P(s)\ge\epsilon,\ \sigma^2(P)\ge\underline\sigma^2\}\).
\par\smallskip\noindent\emph{Member properties.} \texttt{ass:\allowbreak{}fixed-\allowbreak{}face-\allowbreak{}interior}; \texttt{ass:\allowbreak{}nondegenerate-\allowbreak{}variance}.
\end{definition}

\begin{definition}[def:compiler-handle]\label{def:compiler-handle}
\par\smallskip\noindent\emph{Defined object.} \texttt{Support-\allowbreak{}scoped compiler handle}
\par\smallskip\noindent\emph{Construction.} Handle: enrich do-search states with exact prefix domains and signed rational arrays, use the support-saturated matrices \(B_j\) as semantic invariants, and normalize a successful derivation by applying rootwise rational projectors onto \(\operatorname{row}_{\mathbb Q}(B_j)\) before extracting a common denominator.
\par\smallskip\noindent\emph{Inputs.} \texttt{G}; \texttt{S}; \texttt{X}; \texttt{x}; \texttt{Y}; \texttt{y}.
\end{definition}

\section{Main results}

\begin{theorem}[thm:maximal-response-support]\label{thm:maximal-response-support}
Assume \(S\) is realizable. Every \(\theta\in\Theta_S\) satisfies \(\operatorname{supp}(\theta_j)\subseteq T_j(S)\) for all \(j\). Every \(r\in T(S)\) satisfies \(\operatorname{ev}_{\varnothing}(r)\in S\), and \(\operatorname{ev}_{\varnothing}(T(S))=S\). Consequently, every product law strictly positive on each \(T_j(S)\) has factual support exactly \(S\), and \(T(S)\) is the unique largest Cartesian response support among all members of \(\Theta_S\).
\end{theorem}
\begin{proof}
For \(\theta\in\Theta_S\), let
\[
A_j=\operatorname{supp}(\theta_j)
\quad\text{and}\quad
q_j^k(w)=\sum_{r_j\in\mathcal R_j}\theta_j(r_j)b_j^k(w,r_j).
\]
All these sums are finite.  For every prefix \(w\in\prod_{\ell=1}^k\mathcal X_\ell\), acyclicity and the topological order imply the pointwise identity
\[
\mathbf 1\{\operatorname{ev}_{\varnothing}(r)_{\le k}=w\}
=\prod_{j=1}^m b_j^k(w,r_j).                                  \tag{1}
\]
Indeed, the event on the left holds exactly when each response equation at a vertex among \(V_1,\ldots,V_k\) returns the corresponding coordinate of \(w\); the factors on the right partition those equations by district.  Independence of the district roots, encoded by the product law \(\bigotimes_j\theta_j\), and finite distributivity now give
\[
P_\theta(V_{\le k}=w)
=\sum_r\prod_{j=1}^m\theta_j(r_j)b_j^k(w,r_j)
=\prod_{j=1}^m q_j^k(w).                                      \tag{2}
\]

Fix \(i\), let \(V_i\in D_j\), and take \((w,a)\in\mathcal F_i(S)\).  Since \(w\in S^{\le i-1}\), some \(s\in S\) extends \(w\).  Membership \(\theta\in\Theta_S\), through the exact-support condition in \(\mathrm{def:model-class}\), implies
\[
P_\theta(V_{<i}=w)\ge P_\theta(s)>0.
\]
Every factor in (2) is nonnegative, so
\[
q_\ell^{i-1}(w)>0\qquad(\ell=1,\ldots,m).                    \tag{3}
\]
For \(\ell\ne j\), step \(i\) adds no response equation from \(D_\ell\).  Therefore the definition of the prefix indicators gives
\[
b_\ell^i(wa,r_\ell)=b_\ell^{i-1}(w,r_\ell),
\qquad
q_\ell^i(wa)=q_\ell^{i-1}(w)>0.                              \tag{4}
\]
Because \(wa\notin S^{\le i}\) and the factual support is exactly \(S\), equations (2)--(4) yield
\[
0=P_\theta(V_{\le i}=wa)
=q_j^i(wa)\prod_{\ell\ne j}q_\ell^{i-1}(w),
\qquad q_j^i(wa)=0.                                           \tag{5}
\]
The summands \(\theta_j(r_j)b_j^i(wa,r_j)\) defining \(q_j^i(wa)\) are nonnegative.  Hence (5) implies
\[
b_j^i(wa,r_j)=0\qquad\text{for every }r_j\in A_j.
\]
This holds for every first forbidden extension whose new vertex belongs to \(D_j\), so \(\mathrm{def:saturated-local-support}\) gives
\[
A_j\subseteq T_j(S)\qquad(j=1,\ldots,m).                     \tag{6}
\]

Now let \(r\in T(S)\), and put \(v=\operatorname{ev}_{\varnothing}(r)\).  If \(v\notin S\), the convention \(S^{\le0}=\{()\}\) and finiteness give a least \(i\) such that \(v_{\le i}\notin S^{\le i}\).  Then \((v_{<i},v_i)\in\mathcal F_i(S)\).  If \(V_i\in D_j\), evaluation along the realized prefix \(v_{\le i}\) makes every factor in \(b_j^i(v_{\le i},r_j)\) equal to one.  Thus
\[
b_j^i(v_{\le i},r_j)=1,
\]
contradicting \(r_j\in T_j(S)\).  Consequently
\[
\operatorname{ev}_{\varnothing}(T(S))\subseteq S.            \tag{7}
\]

By \(\mathrm{ass:realizable-support}\), choose \(\theta^\star\in\Theta_S\).  For each \(s\in S\), exact support gives \(P_{\theta^\star}(s)>0\).  The finite nonnegative sum defining this probability must therefore contain a tuple \(r\) such that
\[
\prod_j\theta_j^\star(r_j)>0
\quad\text{and}\quad
\operatorname{ev}_{\varnothing}(r)=s.
\]
Such \(r\) lies in \(\prod_j A_j^\star\), which is contained in \(T(S)\) by (6).  Hence \(S\subseteq\operatorname{ev}_{\varnothing}(T(S))\), and (7) proves
\[
\operatorname{ev}_{\varnothing}(T(S))=S.                     \tag{8}
\]

Equation (6) applied to \(\theta^\star\) also shows that every \(T_j(S)\) is nonempty.  Let \(\eta_j\) be any probability mass function strictly positive on every element of \(T_j(S)\) and zero on its complement.  The product law \(\bigotimes_j\eta_j\) is positive at every tuple in \(T(S)\); hence its factual pushforward has support equal to the image in (8), namely \(S\).  It therefore belongs to \(\Theta_S\).  Finally, the joint response support of any \(\theta\in\Theta_S\) is the Cartesian set \(\prod_j A_j\), which is contained in \(T(S)\) by (6), while the constructed law has joint response support exactly \(T(S)\).  Thus \(T(S)\) is a largest Cartesian response support, contains every other such support, and is consequently the unique largest one.
\end{proof}
\par\smallskip\noindent \textit{Justification.} This removes chart ambiguity before any identifying formula is constructed. \textit{Gap.} Zaffalon et al. give local credal constraints and Chen and Darwiche prune zero circuit branches, but neither proves query-independent sufficiency and uniqueness of this support saturation. \textit{Consumer.} The maximal chart is the finite domain on which the tensor identifier and every nonidentification witness are certified.

\begin{theorem}[thm:district-signature-recovery]\label{thm:district-signature-recovery}
For every \(\theta\in\Theta_S\) and every district \(j\), \(B_j\theta_j|_{T_j(S)}=t_j(P_{\theta})\). Moreover, for \(\theta,\theta'\in\Theta_S\), \(P_{\theta}=P_{\theta'}\) if and only if \(B_j\theta_j|_{T_j(S)}=B_j\theta'_j|_{T_j(S)}\) for every \(j\). Every ratio defining \(t_j(P)\) is positive and well defined for every \(P\in\operatorname{relint}(\Delta(S))\).
\end{theorem}
\begin{proof}
Fix \(\theta\in\Theta_S\), and, for every prefix \(w\) of length \(k\), define within this proof
\[
q_j^k(w)=\sum_{r_j\in\mathcal R_j}\theta_j(r_j)b_j^k(w,r_j).
\]
The same acyclic evaluation argument as in (1), followed by the product root law, gives
\[
P_\theta(V_{\le k}=w)=\prod_{\ell=1}^m q_\ell^k(w).          \tag{9}
\]
If \(s\in S\), every prefix event in (9) contains the cell \(s\), and exact support gives \(P_\theta(s)>0\).  Thus
\[
P_\theta(V_{\le k}=s_{\le k})>0,
\qquad
q_\ell^k(s_{\le k})>0                                       \tag{10}
\]
for all \(k\) and \(\ell\), because the factors in (9) are nonnegative.

Suppose \(V_i\in D_j\).  For every \(\ell\ne j\), the definition of \(b_\ell^k\) shows that its value does not change from step \(i-1\) to step \(i\).  Cancelling the positive common factors in (9), justified by (10), yields
\[
\frac{P_\theta(V_{\le i}=s_{\le i})}
     {P_\theta(V_{<i}=s_{<i})}
=\frac{q_j^i(s_{\le i})}{q_j^{i-1}(s_{<i})}.                 \tag{11}
\]
The sequence \(q_j^k(s_{\le k})\) changes only at indices \(k\) for which \(V_k\in D_j\), and \(q_j^0(())=1\) because the defining indicator is an empty product.  Multiplying (11) over all such indices therefore telescopes:
\[
\prod_{i:\,V_i\in D_j}
\frac{P_\theta(V_{\le i}=s_{\le i})}
     {P_\theta(V_{<i}=s_{<i})}
=q_j^n(s)
=\sum_{r_j\in\mathcal R_j}\theta_j(r_j)b_j^n(s,r_j).        \tag{12}
\]
By \(\mathrm{thm:maximal-response-support}\), \(\theta_j\) vanishes outside \(T_j(S)\).  The right side of (12) is consequently the \(s\)-coordinate of \(B_j\theta_j|_{T_j(S)}\), and the left side is \(t_j(P_\theta)[s]\).  Their normalization coordinates also agree because
\[
(B_j\theta_j|_{T_j(S)})[0]
=\sum_{r_j\in T_j(S)}\theta_j(r_j)=1=t_j(P_\theta)[0].
\]
It follows that
\[
B_j\theta_j|_{T_j(S)}=t_j(P_\theta).                          \tag{13}
\]

If \(P_\theta=P_{\theta'}\), positivity in (10) makes every ratio defining \(t_j\) well defined, so (13) gives equality of the district signatures.  Conversely, for every \(s\in S\), equations (9) and (12) give
\[
P_\theta(s)=\prod_{j=1}^m q_j^n(s)
=\prod_{j=1}^m(B_j\theta_j|_{T_j(S)})[s].                    \tag{14}
\]
Thus equality of every district signature for \(\theta\) and \(\theta'\) implies equality of their probabilities at each \(s\in S\).  Both laws vanish on \(\mathcal X_V\setminus S\) by membership in \(\Theta_S\), so they are equal everywhere.  This proves the claimed equivalence.

Finally, let \(P\in\operatorname{relint}(\Delta(S))\).  For \(s\in S\) and every \(i\),
\[
P(V_{<i}=s_{<i})
=\sum_{u\in S:\,u_{<i}=s_{<i}}P(u)
\ge P(s)>0.                                                    \tag{15}
\]
The numerator \(P(V_{\le i}=s_{\le i})\) is likewise at least \(P(s)>0\).  Hence every ratio in every \(t_j(P)[s]\) is defined and strictly positive throughout the relative interior, including when \(P\) is not itself compatible with an SCM.
\end{proof}
\par\smallskip\noindent \textit{Justification.} Positive-prefix telescoping converts the multilinear observational map into separate identified linear signatures. \textit{Gap.} Hwang et al. provide zero-aware \(Q\)-factor cancellation and Zaffalon et al. provide rootwise credal constraints, but the exact-support fiber equivalence on the maximal chart is not stated in either source. \textit{Consumer.} It reduces same-law comparison from a nonlinear product fiber to finite rational linear algebra district by district.

\begin{theorem}[thm:rational-functional]\label{thm:rational-functional}
If \(h_{x,y}|_{T(S)}\in\operatorname{row}_{\mathbb Q}(B)\), then a rational vector \(\lambda\) satisfying \(\lambda B=h_{x,y}|_{T(S)}\) exists and \(\tau_{x,y}(\theta)=F_{\lambda}(P_{\theta})\) for every \(\theta\in\Theta_S\). The expression has a rational-coefficient numerator over a denominator dividing \(D_S(P)\), and \(D_S(P)>0\) throughout \(\operatorname{relint}(\Delta(S))\).
\end{theorem}
\begin{proof}
For \(\theta\in\Theta_S\), write \(\vartheta_j=\theta_j|_{T_j(S)}\) as a column vector.  The normalization coordinate of the district-signature identity gives
\[
\sum_{r_j\in T_j(S)}\vartheta_j(r_j)=1.
\]
Since \(\theta_j\) is a probability mass function, it follows that it vanishes outside \(T_j(S)\).  The definitions of \(h_{x,y}\), \(\tau_{x,y}\), and the Kronecker matrix therefore imply
\[
\tau_{x,y}(\theta)
=h_{x,y}|_{T(S)}\left(\bigotimes_{j=1}^m\vartheta_j\right),
\qquad
B\left(\bigotimes_{j=1}^m\vartheta_j\right)
=\bigotimes_{j=1}^m(B_j\vartheta_j).                          \tag{16}
\]
The entries of \(B\) and \(h_{x,y}|_{T(S)}\) are integers.  The hypothesis
\(h_{x,y}|_{T(S)}\in\operatorname{row}_{\mathbb Q}(B)\) means exactly that the finite rational linear system
\[
\lambda B=h_{x,y}|_{T(S)}                                    \tag{17}
\]
has a solution \(\lambda\) with rational coordinates.  Multiplying (17) by the product mass vector in (16), using the Kronecker identity in (16), and then invoking \(\mathrm{thm:district-signature-recovery}\), gives
\[
\begin{aligned}
\tau_{x,y}(\theta)
&=\lambda B\left(\bigotimes_j\vartheta_j\right)\\
&=\sum_{(a_1,\ldots,a_m)\in(\{0\}\cup S)^m}
  \lambda[a_1,\ldots,a_m]\prod_{j=1}^m(B_j\vartheta_j)[a_j]\\
&=\sum_{(a_1,\ldots,a_m)\in(\{0\}\cup S)^m}
  \lambda[a_1,\ldots,a_m]\prod_{j=1}^m t_j(P_\theta)[a_j]\\
&=F_\lambda(P_\theta).                                       \tag{18}
\end{aligned}
\]
Thus the displayed observed-law functional identifies the causal target on \(\Theta_S\); the equality was derived rather than imposed as a definition.

It remains to establish the claimed global chart.  Each marginal probability in \(t_j(P)\) is a sum of cells \(P(s)\), hence a polynomial with integer coefficients.  In a fixed product \(\prod_jt_j(P)[a_j]\), the denominator indexed by a vertex \(V_i\) occurs only in the factor belonging to the unique district containing \(V_i\).  When \(a_j=s\in S\), that denominator is
\[
P(V_{<i}=s_{<i}),
\qquad s_{<i}\in S^{\le i-1};
\]
when \(a_j=0\), the factor is one and contributes no denominator.  Hence no factor indexed by \((i,w)\) occurs more than once, and the denominator of every product term divides
\[
D_S(P)=\prod_{i=1}^n\prod_{w\in S^{\le i-1}}P(V_{<i}=w).
\]
Because \(\lambda\) is rational and the sum is finite, bringing (18) to this common denominator gives
\[
F_\lambda(P)=\frac{N_\lambda(P)}{D_S(P)}                     \tag{19}
\]
on the relative interior for a polynomial \(N_\lambda\) with rational coefficients; cancellation can only replace \(D_S\) by a divisor.

For each \(w\in S^{\le i-1}\), choose an extending cell \(s\in S\).  If \(P\in\operatorname{relint}(\Delta(S))\), then
\[
P(V_{<i}=w)\ge P(s)>0.                                       \tag{20}
\]
Every nonempty-prefix factor of \(D_S(P)\) is therefore positive, while the empty-prefix factor is one.  Thus \(D_S(P)>0\) throughout \(\operatorname{relint}(\Delta(S))\), which proves the denominator assertion.
\end{proof}
\par\smallskip\noindent \textit{Justification.} Success must yield one evaluable formula valid everywhere on the face, not a numerical singleton test. \textit{Gap.} Arithmetic-circuit invariant cuts and credal singleton computations are instance certificates; the common denominator and face-wide rational identity are the additional output here. \textit{Consumer.} The functional can be evaluated, differentiated, audited symbolically, and passed to a support-scoped identification engine.

\begin{theorem}[thm:rational-interior-twins]\label{thm:rational-interior-twins}
If \(h_{x,y}|_{T(S)}\notin\operatorname{row}_{\mathbb Q}(B)\), then there exist \(\theta^+,\theta^-\in\Theta_S\) with rational coordinates, \(\theta_j^+(r_j)>0\) and \(\theta_j^-(r_j)>0\) for every \(j\) and every \(r_j\in T_j(S)\), such that \(P_{\theta^+}=P_{\theta^-}\) and \(\tau_{x,y}(\theta^+)\ne\tau_{x,y}(\theta^-)\). The pair may be chosen to differ in exactly one district law.
\end{theorem}
\begin{proof}
Write \(E_j=\mathbb Q^{T_j(S)}\) and \(L_j=\operatorname{row}_{\mathbb Q}(B_j)\subseteq E_j\).  Realizability and the maximal-response-support theorem imply that every \(T_j(S)\) is nonempty: choose any \(\bar\theta\in\Theta_S\); then \(\operatorname{supp}(\bar\theta_j)\) is a nonempty subset of \(T_j(S)\).  We identify \(\mathbb Q^{T(S)}\) with \(\bigotimes_{j=1}^m E_j\) by the standard coordinate basis.

For each \(j\), choose a basis of \(L_j\) and extend it to a basis of \(E_j\).  The resulting tensor-product basis proves
\[
\bigcap_{j=1}^m
\left(E_1\otimes\cdots\otimes E_{j-1}\otimes L_j
\otimes E_{j+1}\otimes\cdots\otimes E_m\right)
=\bigotimes_{j=1}^m L_j.                                      \tag{1}
\]
Indeed, membership in the \(j\)-th space on the left forces to zero every tensor-basis coordinate whose \(j\)-th factor belongs to the chosen complement of \(L_j\); imposing this for every \(j\) leaves precisely the coordinates with all factors in their corresponding \(L_j\).  Since the rows of a Kronecker product are exactly the tensor products of rows of its factors, the definition of \(B\) gives
\[
\operatorname{row}_{\mathbb Q}(B)=\bigotimes_{j=1}^mL_j.     \tag{2}
\]

Set \(h=h_{x,y}|_{T(S)}\).  Under the hypothesis \(h\notin\operatorname{row}_{\mathbb Q}(B)\), equations (1)--(2) give an index \(j\) such that
\[
h\notin E_1\otimes\cdots\otimes E_{j-1}\otimes L_j
\otimes E_{j+1}\otimes\cdots\otimes E_m.                    \tag{3}
\]
To turn (3) into a rational separator, expand \(h\) in a rational basis of \(\bigotimes_{k\ne j}E_k\).  At least one associated coefficient vector in \(E_j\) is outside \(L_j\).  Rational Gaussian elimination therefore supplies a vector
\[
v\in L_j^\perp\cap\mathbb Q^{T_j(S)}=\ker(B_j)
\]
whose contraction with that coefficient vector is nonzero.  Equivalently, the rational tensor
\[
g(r_{-j})=
\sum_{r_j\in T_j(S)}h(r_j,r_{-j})v(r_j)                       \tag{4}
\]
is not identically zero on \(\prod_{k\ne j}T_k(S)\).  The equality \(L_j^\perp=\ker(B_j)\) follows directly from the coordinate pairing: \(v\) annihilates every row of \(B_j\) exactly when \(B_jv=0\).  Because the row \(B_j[0,\cdot]\) is the all-ones row, we also have
\[
\sum_{r_j\in T_j(S)}v(r_j)=0.                                \tag{5}
\]

For \(k\ne j\), let \(d_k=|T_k(S)|\), let \(u_k\) be the uniform probability vector on \(T_k(S)\), and, for \(a\in T_k(S)\), define
\[
\eta_{k,a}=\frac{u_k+e_a}{2}.
\]
Every \(\eta_{k,a}\) is a strictly positive rational probability vector.  The matrix having these vectors as columns is
\[
\frac12\left(I+u_k\mathbf 1^{\mathsf T}\right),
\]
and \(I+u_k\mathbf 1^{\mathsf T}\) acts as multiplication by two on \(\operatorname{span}\{u_k\}\) and as the identity on \(\ker(\mathbf 1^{\mathsf T})\).  These two spaces form a direct sum because \(\mathbf 1^{\mathsf T}u_k=1\).  Hence its determinant is two, so the determinant of the displayed matrix is
\[
2^{-d_k}\left(1+\mathbf 1^{\mathsf T}u_k\right)=2^{1-d_k}>0.
\]
Thus the vectors \(\{\eta_{k,a}:a\in T_k(S)\}\) span \(E_k\), and their pure tensor products span \(\bigotimes_{k\ne j}E_k\).  Since \(g\ne0\), some choice \((a_k:k\ne j)\) satisfies
\[
c=\sum_{r_{-j}}g(r_{-j})
       \prod_{k\ne j}\eta_{k,a_k}(r_k)\ne0.                 \tag{6}
\]
When \(m=1\), the empty product is one and (6) is simply the nonzero scalar contraction in (4).

Let \(u_j\) be uniform on \(T_j(S)\).  Finiteness, strict positivity of \(u_j\), and rationality of \(v\) allow a rational \(\delta>0\) such that
\[
\delta |v(r_j)|<u_j(r_j)
\quad\text{for every }r_j\text{ with }v(r_j)\ne0.            \tag{7}
\]
Define probability vectors on the full response spaces by
\[
\theta_j^+=u_j+\delta v,
\qquad
\theta_j^-=u_j-\delta v,
\qquad
\theta_k^+=\theta_k^-=\eta_{k,a_k}\quad(k\ne j),             \tag{8}
\]
and set every coordinate outside \(T_k(S)\) equal to zero.  Equation (5) proves normalization, while (7) proves strict positivity on every \(T_k(S)\).  All coordinates in (8) are rational.  The maximal-response-support theorem now applies to each product law in (8), proving
\[
\theta^+,\theta^-\in\Theta_S.                                \tag{9}
\]
Moreover,
\[
B_j\theta_j^+-B_j\theta_j^-=2\delta B_jv=0,                  \tag{10}
\]
and the corresponding equality for \(k\ne j\) is immediate from (8).  The district-signature-recovery theorem and (9)--(10) therefore give
\[
P_{\theta^+}=P_{\theta^-}.                                   \tag{11}
\]
Finally, the definition of \(\tau_{x,y}\), followed by (4), (6), and (8), yields
\[
\begin{aligned}
\tau_{x,y}(\theta^+)-\tau_{x,y}(\theta^-)
&=\sum_{r_{-j}}\sum_{r_j}h(r_j,r_{-j})
  \bigl(2\delta v(r_j)\bigr)
  \prod_{k\ne j}\eta_{k,a_k}(r_k)\\
&=2\delta c\ne0.                                             \tag{12}
\end{aligned}
\]
In particular \(v\ne0\), so (8) differs in exactly the \(j\)-th district law.  This proves all assertions of the theorem.

The same constructed pair also gives the directed ambient-class consequence.  If \(\Theta_S\subseteq\mathcal M_S\) and the embedding into \(\mathcal M_S\) preserves the observational and \(\operatorname{do}(X=x)\) semantics, then (9), (11), and (12) hold for the same two elements regarded as members of \(\mathcal M_S\).  Hence nonmembership in the tensor row span refutes uniform identification in every such ambient class.  This uses only inclusion, so it supplies no converse transfer of the success direction.
\end{proof}
\par\smallskip\noindent \textit{Justification.} The separator construction preserves district independence, strict positivity on the saturated chart, rationality, exact support, and the exact observational law inside the canonical quasi-Markovian class \(\Theta_S\); inclusion then carries the same witnesses into every semantics-preserving supermodel class. \textit{Gap.} Classical hedges need not preserve a prescribed zero face, while generic polynomial or credal optimization need not emit exact symbolic interior twins. The new ambient conclusion is deliberately converse-only. \textit{Consumer.} Every failed tensor test yields reproducible finite SCM twins and therefore a nonidentification witness both on \(\Theta_S\) and in every admissible ambient class containing it.

\begin{theorem}[thm:uniform-id-iff-tensor-span]\label{thm:uniform-id-iff-tensor-span}
The query \(\tau_{x,y}\) is uniformly identified on the exact support face, meaning \(\tau_{x,y}(\theta)=\tau_{x,y}(\theta')\) for all \(\theta,\theta'\in\Theta_S\) with \(P_{\theta}=P_{\theta'}\), if and only if \(h_{x,y}|_{T(S)}\in\operatorname{row}_{\mathbb Q}(B)=\bigotimes_{j=1}^m\operatorname{row}_{\mathbb Q}(B_j)\).
\end{theorem}
\begin{proof}
By the definition \(B=\bigotimes_{j=1}^mB_j\), every row of \(B\) is a tensor product of one row from each \(B_j\), and every such tensor product occurs as a row.  Bilinearity of the tensor product therefore gives
\[
\operatorname{row}_{\mathbb Q}(B)
=\bigotimes_{j=1}^m\operatorname{row}_{\mathbb Q}(B_j).       \tag{13}
\]

Suppose first that \(h_{x,y}|_{T(S)}\in\operatorname{row}_{\mathbb Q}(B)\).  The rational-functional theorem supplies \(\lambda\) and a single observed-law map \(F_\lambda\) satisfying
\[
\tau_{x,y}(\theta)=F_\lambda(P_\theta)
\quad\text{for every }\theta\in\Theta_S.                    \tag{14}
\]
Thus, for \(\theta,\theta'\in\Theta_S\), the equality \(P_\theta=P_{\theta'}\) implies, by applying the same map to both sides of (14), that \(\tau_{x,y}(\theta)=\tau_{x,y}(\theta')\).  This is uniform identification on the canonical quasi-Markovian class \(\Theta_S\).

Conversely, if \(h_{x,y}|_{T(S)}\notin\operatorname{row}_{\mathbb Q}(B)\), the rational-interior-twins theorem produces \(\theta^+,\theta^-\in\Theta_S\) for which
\[
P_{\theta^+}=P_{\theta^-}
\quad\text{but}\quad
\tau_{x,y}(\theta^+)\ne\tau_{x,y}(\theta^-).                \tag{15}
\]
Equation (15) contradicts the defining same-law implication for uniform identification on \(\Theta_S\).  The two row-space cases are exhaustive, and (13) proves the displayed tensor equality, establishing the stated equivalence.

For completeness about class scope, the ambient-supermodel converse-transfer proposition applies to (15): for every finite-SCM class \(\mathcal M_S\) containing \(\Theta_S\) and preserving its observational and intervention semantics, uniform identification over \(\mathcal M_S\) implies \(h_{x,y}|_{T(S)}\in\operatorname{row}_{\mathbb Q}(B)\).  The argument above proves the reverse implication only on \(\Theta_S\); it does not assert that tensor membership identifies the query over a broader class with overlapping latent-parent neighborhoods.
\end{proof}
\par\smallskip\noindent \textit{Justification.} This is the complete formula-or-twins characterization on the canonical quasi-Markovian class; its declared ambient corollary transfers only the necessary direction to supermodel classes. \textit{Gap.} Nearby positivity, invariant, and optimization results do not give this support-uniform tensor equivalence on \(\Theta_S\); the theorem does not claim sufficiency for broader overlapping-latent classes. \textit{Consumer.} It decides identification on \(\Theta_S\) and, through the ambient converse, supplies a necessary test for any semantics-preserving finite-SCM supermodel class.

\begin{proposition}[prop:destructive-controls]\label{prop:destructive-controls}
The tensor criterion has the following reductions. For \((G_{\mathrm{bow}},S_{\mathrm{diag}})\), it rejects identification of \(P_{\operatorname{do}(X=1)}(Y=1)\). For \((G_{\mathrm{IV}},S_{\mathrm{IV5}})\), it identifies \(P_{\operatorname{do}(X=0)}(Y=1)=P(X=0,Y=1\mid Z=1)\). When \(S=\mathcal X_V\), every classically identified query is accepted and the rational functional agrees with the classical identifying functional on positive laws.
\end{proposition}
\begin{proof}
For \((G_{\mathrm{bow}},S_{\mathrm{diag}})\), let \(U\) be a fair binary root for the sole district.  Consider the two finite acyclic SCMs
\[
\mathcal M_1:\quad X=U,\quad Y=U,
\qquad
\mathcal M_2:\quad X=U,\quad Y=X.                            \tag{21}
\]
Ignoring a displayed parent is permitted by the graph, so both SCMs are compatible with \(G_{\mathrm{bow}}\).  Encoding each value of \(U\) by its deterministic response tables places both models in the canonical response class.  In both cases the factual law assigns mass \(1/2\) to each of \((0,0)\) and \((1,1)\), and zero elsewhere, so its exact support is \(S_{\mathrm{diag}}\).  Under \(\operatorname{do}(X=1)\), however,
\[
P_{\mathcal M_1}(Y=1\mid\operatorname{do}(X=1))=\frac12,
\qquad
P_{\mathcal M_2}(Y=1\mid\operatorname{do}(X=1))=1.           \tag{22}
\]
Thus the query is not uniformly identified on the exact-support class.  The reverse implication in \(\mathrm{thm:uniform-id-iff-tensor-span}\) says that tensor membership would imply uniform identification; its contrapositive therefore makes the tensor criterion reject this query.  Equation (22) also proves directly that the factual relation \(Y=X\) on the diagonal support cannot be substituted after intervention.

For \((G_{\mathrm{IV}},S_{\mathrm{IV5}})\), the districts are \(\{Z\}\) and \(\{X,Y\}\).  Write a downstream deterministic response type as
\[
(x_0,x_1,y_0,y_1),
\]
where \(x_z\) is the response of \(X\) to \(Z=z\), and \(y_x\) is the response of \(Y\) to \(X=x\).  The first-forbidden-prefix conditions can be read off exactly.  At \(z=0\), the factual support permits only \((x,y)=(0,0)\) and \((1,0)\), so it imposes \(y_0=0\) when \(x_0=0\) and \(y_1=0\) when \(x_0=1\).  At \(z=1\), it permits \((0,0),(0,1),(1,1)\), so it imposes no restriction on \(y_0\) when \(x_1=0\), and imposes \(y_1=1\) when \(x_1=1\).  Splitting on the four values of \((x_0,x_1)\) gives precisely
\[
\begin{array}{cc}
(0,0,0,0),&(0,0,0,1),\\
(0,1,0,1),&(1,0,0,0),\\
(1,0,1,0),
\end{array}                                                    \tag{23}
\]
with \((x_0,x_1)=(1,1)\) impossible because it would require both \(y_1=0\) and \(y_1=1\).  Thus (23) is exactly the downstream saturated local response support.

Direct inspection of all five rows of (23) proves the pointwise identity
\[
\mathbf 1\{y_0=1\}
=\mathbf 1\{x_1=0,\ y_{x_1}=1\}.                            \tag{24}
\]
The left side is the response indicator for \(Y=1\) under \(\operatorname{do}(X=0)\).  The right side is the response indicator for \((X,Y)=(0,1)\) under \(\operatorname{do}(Z=1)\).  The root for \(Z\) is independent of the downstream district root under the product parameterization.  Moreover, exact support and the presence of cells with \(Z=1\) in \(S_{\mathrm{IV5}}\) imply \(P(Z=1)>0\).  Consequently factual conditioning on \(Z=1\) leaves the downstream root law unchanged, and taking expectations in (24) gives
\[
\begin{aligned}
P_{\operatorname{do}(X=0)}(Y=1)
&=P_{\operatorname{do}(Z=1)}(X=0,Y=1)\\
&=P(X=0,Y=1\mid Z=1).                                        \tag{25}
\end{aligned}
\]
This equality holds for every member of the exact-support class.  The forward direction of \(\mathrm{thm:uniform-id-iff-tensor-span}\) therefore makes the tensor criterion accept the query, and \(\mathrm{thm:rational-functional}\) compiles the accepted row certificate.  The last expression in (25) is well defined because its denominator \(P(Z=1)\) is strictly positive.

Finally suppose \(S=\mathcal X_V\), and let the query be classically identified for \(G\).  Classical identification means that two graph-compatible SCMs with the same observational law have the same value of the query.  In particular this implication holds on the canonical independent-district-root subclass \(\Theta_S\), so the query is uniformly identified there.  By \(\mathrm{thm:uniform-id-iff-tensor-span}\), the tensor criterion accepts it.  Conversely, every finite SCM in the stated independent-root-per-district model can be put in canonical response form: map each value of a district root to the tuple of deterministic response tables that it induces, and push the root law forward under this map.  Independence makes the resulting response law a product \(\bigotimes_j\theta_j\), and the construction preserves both factual and interventional evaluations.  Hence every positive graph-compatible observational law in the stated model has a representative \(\theta\in\Theta_S\).  At such a law, the classical identifying functional and the rational functional from \(\mathrm{thm:rational-functional}\) both equal the same interventional probability, so they agree.  This comparison is on positive graph-compatible laws; no equality away from the observational model is needed.
\end{proof}
\par\smallskip\noindent \textit{Justification.} These controls prevent factual equalities or unsupported cancellations from masquerading as intervention identities. \textit{Gap.} The five-cell identity and diagonal-bow rejection are not supplied by formula-level positivity checks as a paired completeness diagnostic. \textit{Consumer.} They provide minimal regression tests for any implementation and recover the familiar full-support regime.

\begin{theorem}[thm:face-multinomial-inference]\label{thm:face-multinomial-inference}
Fix a successful certificate \(\lambda\). In the saturated multinomial model on \(S\), \(\phi_P\) is the efficient influence function of \(F_{\lambda}(P)\). Under independent multinomial sampling, uniformly over \(P\in\mathcal P_{S,\epsilon,\underline\sigma}\), \(\sqrt N\{\widehat\tau_N-F_{\lambda}(P)\}=N^{-1/2}\sum_{k=1}^N\phi_P(O_k)+o_P(1)\), the studentized statistic converges uniformly to \(\mathcal N(0,1)\), and the plug-in variance \(\sum_{s\in S}\widehat P_N(s)\phi_{\widehat P_N}(s)^2\) is uniformly consistent. On the compatible submodel this gives Wald inference for \(\tau_{x,y}\).
\end{theorem}
\begin{proof}
Let \(d=|S|\), identify a mass function on \(S\) with a vector in \(\mathbb R^d\), and put
\[
K_\epsilon=\left\{p\in\Delta(S):\min_{s\in S}p(s)\ge\epsilon\right\}.
\]
If \(w\in S^{\le i-1}\), choose \(s\in S\) extending \(w\).  Then, for every \(p\in K_\epsilon\),
\[
p(V_{<i}=w)\ge p(s)\ge\epsilon.                              \tag{26}
\]
The rational representation in \(\mathrm{thm:rational-functional}\) and (26) imply that \(F_\lambda\) is twice continuously differentiable on a relative neighborhood of \(K_\epsilon\), with its first and second partial derivatives uniformly bounded there.  For example, the same conclusion holds on the compact set where every cell mass is at least \(\epsilon/2\), because all supported-prefix denominators are then at least \(\epsilon/2\).

Write
\[
g_P(s)=\frac{\partial F_\lambda(P)}{\partial P(s)}.
\]
By \(\mathrm{def:influence-function}\),
\[
\phi_P(s)=g_P(s)-\sum_{u\in S}P(u)g_P(u),
\qquad
\sum_{s\in S}P(s)\phi_P(s)=0.                               \tag{27}
\]
Let \(a:S\to\mathbb R\) be an arbitrary score in the saturated multinomial tangent space, so \(\sum_sP(s)a(s)=0\).  Since \(P(s)>0\) on the relative interior, \(P_t(s)=P(s)\{1+t a(s)\}\) is a valid submodel for all sufficiently small \(|t|\).  The chain rule and (27) give
\[
\left.\frac{d}{dt}F_\lambda(P_t)\right|_{t=0}
=\sum_sg_P(s)P(s)a(s)
=\sum_s\phi_P(s)P(s)a(s).                                    \tag{28}
\]
Every mean-zero function on the finite set \(S\) occurs as such a tangent score.  Therefore \(\phi_P\) is the unique mean-zero representer of every pathwise derivative in the saturated multinomial tangent space.  It is consequently the efficient influence function, and its variance is \(\sigma^2(P)\) by definition.

Define
\[
E_N=\left\{\|\widehat P_N-P\|_\infty\le\epsilon/2\right\}.
\]
Under \(\mathrm{ass:iid-multinomial-sampling}\), the multinomial variance formula, Chebyshev's inequality, and a union bound give
\[
\sup_{P\in K_\epsilon}P(E_N^c)
\le \frac{4d}{N\epsilon^2}\longrightarrow0.                 \tag{29}
\]
On \(E_N\), every empirical cell mass is at least \(\epsilon/2\), so \(D_S(\widehat P_N)>0\) by (26), and \(\mathrm{def:plugin-estimator}\) gives \(\widehat\tau_N=F_\lambda(\widehat P_N)\).  The uniform second-derivative bound and Taylor's theorem yield
\[
F_\lambda(\widehat P_N)-F_\lambda(P)
=\sum_{s\in S}g_P(s)\{\widehat P_N(s)-P(s)\}+R_N,
\qquad
|R_N|\le C\|\widehat P_N-P\|_2^2,                            \tag{30}
\]
where \(C<\infty\) is independent of \(P\in K_\epsilon\).  Because the empirical and true masses both sum to one, (27) turns the linear term into
\[
\sum_s\phi_P(s)\{\widehat P_N(s)-P(s)\}
=\frac1N\sum_{k=1}^N\phi_P(O_k).                             \tag{31}
\]
The multinomial variance formula also gives
\[
\sup_{P\in K_\epsilon}N\,\mathbb E_P
   \|\widehat P_N-P\|_2^2
=\sup_{P\in K_\epsilon}\sum_{s\in S}P(s)\{1-P(s)\}
\le d.                                                        \tag{32}
\]
For every \(\delta>0\), Markov's inequality and (32) imply
\[
\sup_{P\in K_\epsilon}
P\!\left(\sqrt N\|\widehat P_N-P\|_2^2>\delta\right)
\le\frac{d}{\delta\sqrt N}\longrightarrow0.                \tag{33}
\]
Combining (29)--(33) proves, uniformly over the smaller class \(\mathcal P_{S,\epsilon,\underline\sigma}\subseteq K_\epsilon\),
\[
\sqrt N\{\widehat\tau_N-F_\lambda(P)\}
=\frac1{\sqrt N}\sum_{k=1}^N\phi_P(O_k)+o_P(1).             \tag{34}
\]

Uniform boundedness of the first derivatives on \(K_\epsilon\), (27), and the condition \(\sigma(P)\ge\underline\sigma>0\) imply that a constant \(M<\infty\) satisfies
\[
\left|\frac{\phi_P(s)}{\sigma(P)}\right|\le M
\quad\text{for all }s\in S,
\quad P\in\mathcal P_{S,\epsilon,\underline\sigma}.       \tag{35}
\]
For any sequence \(P_N\) in this class, set
\[
W_{N,k}=\frac{\phi_{P_N}(O_k)}{\sqrt N\,\sigma(P_N)},
\qquad k=1,\ldots,N.
\]
The variables in each row are independent and mean zero by (27), and their total variance is one.  For any \(\delta>0\), (35) implies \(|W_{N,k}|\le M/\sqrt N\le\delta\) for all sufficiently large \(N\); hence the Lindeberg sum is eventually zero.  The cited \(\mathrm{lem:lindeberg-feller-triangular-array}\) gives
\[
\frac1{\sqrt N\,\sigma(P_N)}
\sum_{k=1}^N\phi_{P_N}(O_k)
\rightsquigarrow\mathcal N(0,1).                             \tag{36}
\]
If convergence in (36) were not uniform over \(\mathcal P_{S,\epsilon,\underline\sigma}\), one could select a violating sequence \(P_N\), contradicting (36).  Thus the convergence is uniform.

On the relative neighborhood used above, define
\[
v(p)=\sum_{s\in S}p(s)\phi_p(s)^2.
\]
The positive denominator bound and continuity of the rational derivatives make \(v\) uniformly continuous on the compact set where all cell masses are at least \(\epsilon/2\).  Equations (29) and (32) imply \(\widehat P_N-P=o_P(1)\) uniformly, and therefore
\[
\widehat v_N
:=\sum_{s\in S}\widehat P_N(s)\phi_{\widehat P_N}(s)^2
=\sigma^2(P)+o_P(1)                                          \tag{37}
\]
uniformly on the event \(D_S(\widehat P_N)>0\).  Assigning any fixed value to \(\widehat v_N\) on its complement does not affect (37), because that complement is contained in \(E_N^c\), whose probability vanishes uniformly by (29).  Since \(\sigma^2(P)\ge\underline\sigma^2\), equations (34), (36), and (37), together with Slutsky's argument applied to arbitrary sequences \(P_N\), prove uniform convergence of
\[
\frac{\sqrt N\{\widehat\tau_N-F_\lambda(P)\}}
     {\widehat v_N^{1/2}}
\]
to \(\mathcal N(0,1)\), as well as uniform consistency of the displayed plug-in variance.  Finally, on the compatible submodel, \(\mathrm{thm:rational-functional}\) gives
\[
F_\lambda(P_\theta)=\tau_{x,y}(\theta)
=P_\theta(Y=y\mid\operatorname{do}(X=x)).
\]
Thus the expansion, studentization, and Wald intervals concern the causal target on that submodel, while \(F_\lambda(P)\) remains the distinct observed-law functional on the full saturated face.
\end{proof}
\par\smallskip\noindent \textit{Justification.} A complete symbolic identifier becomes statistically usable once its rational output has face-tangent inference. \textit{Gap.} Bhattacharya, Nabi, and Shpitser derive graph-native influence functions for important functional classes, but not for arbitrary rational outputs created by exact-support identification on a lower-dimensional multinomial face. \textit{Consumer.} Analysts with known structural zeros obtain standard errors and confidence intervals without replacing the deterministic intervention target.

\begin{openendedquestion}[oeq:support-scoped-compiler]\label{oeq:support-scoped-compiler}
Does the support-scoped compiler handle yield a finite do-search calculus that is sound and complete for \(\tau_{x,y}\) over \(\Theta_S\), permits signed intermediate expressions without invalid cancellation, and normalizes every successful derivation to a bounded certificate \(\lambda B=h_{x,y}|_{T(S)}\) whose denominators divide \(D_S(P)\)? Establish termination, rule-by-rule soundness, completeness relative to the tensor criterion, and an explicit certificate-size bound, or exhibit the first finite counterexample.
\end{openendedquestion}
\par\smallskip\noindent \textit{Justification.} The unresolved content is whether product-preserving signed search rules can realize the complete semantic criterion without materializing the full response tensor. \textit{Gap.} Do-search inherits completeness only from its underlying calculus, and positivity-aware work does not compile structural-zero gains into tensor-normal certificates. \textit{Consumer.} A positive answer gives a support-aware extension of do-search with independently checkable success and failure receipts.

\begin{lemma}[lem:lindeberg-feller-triangular-array]\label{lem:lindeberg-feller-triangular-array}
Let \(W_{q,1},\ldots,W_{q,k_q}\) be a triangular array of independent real random variables with mean zero.  If \(\sum_{i=1}^{k_q}\operatorname{Var}(W_{q,i})\to1\) and, for every \(\delta>0\),
\[
\sum_{i=1}^{k_q}\mathbb E\!\left[W_{q,i}^2\mathbf 1\{|W_{q,i}|>\delta\}\right]\longrightarrow0,
\]
then \(\sum_{i=1}^{k_q}W_{q,i}\rightsquigarrow\mathcal N(0,1)\).
\end{lemma}

\begin{proposition}[prop:ambient-supermodel-converse-transfer]\label{prop:ambient-supermodel-converse-transfer}
Let \(\mathcal M_S\) be an ambient finite-SCM class with \(\Theta_S\subseteq\mathcal M_S\) such that the observational law and the \(\operatorname{do}(X=x)\) target of every embedded \(\theta\in\Theta_S\) are respectively \(P_\theta\) and \(\tau_{x,y}(\theta)\).  If \(h_{x,y}|_{T(S)}\notin\operatorname{row}_{\mathbb Q}(B)\), then there exist \(\theta^+,\theta^-\in\Theta_S\subseteq\mathcal M_S\), rational and strictly positive on every \(T_j(S)\), that differ in exactly one district law and satisfy \(P_{\theta^+}=P_{\theta^-}\) but \(\tau_{x,y}(\theta^+)\ne\tau_{x,y}(\theta^-)\).  Consequently the query is not uniformly identified over \(\mathcal M_S\).  Equivalently, uniform identification over any such \(\mathcal M_S\) implies \(h_{x,y}|_{T(S)}\in\operatorname{row}_{\mathbb Q}(B)\).
\end{proposition}
\begin{proof}
Assume tensor nonmembership.  The rational-interior-twins theorem supplies \(\theta^+,\theta^-\in\Theta_S\) with rational coordinates, strict positivity on each \(T_j(S)\), equality of observational laws, unequal target values, and a difference in exactly one district law.  The inclusion \(\Theta_S\subseteq\mathcal M_S\) places these same two models in \(\mathcal M_S\).  Preservation of semantics under the embedding leaves both the equality \(P_{\theta^+}=P_{\theta^-}\) and the inequality \(\tau_{x,y}(\theta^+)\ne\tau_{x,y}(\theta^-)\) unchanged.  They therefore witness failure of the same-law implication defining uniform identification over \(\mathcal M_S\).  Taking the contrapositive gives the final implication.  No step proves that tensor membership is sufficient on \(\mathcal M_S\), so the success direction remains scoped to \(\Theta_S\).
\end{proof}
\par\smallskip\noindent \textit{Justification.} The negative certificate is inherited by every supermodel class containing the canonical class because its two witnesses remain available with unchanged semantics. \textit{Gap.} This transfer is asymmetric: inclusion propagates counterexamples upward but does not propagate an identifying functional from the subclass to the supermodel class. \textit{Consumer.} It makes the tensor condition a necessary condition for identification in broader finite-SCM classes without overstating sufficiency there.

\section*{Technical internal limitation (diagnostic only)}

The iff and rational success certificate are proved for the canonical quasi-Markovian class \(\Theta_S\) with one independent finite root per bidirected-clique district. Inclusion of \(\Theta_S\) in a broader semantics-preserving class transfers the negative twins, hence necessity of tensor membership, but cannot transfer sufficiency because the ambient class may contain additional same-law models. The semantic tensor criterion also permits signed rational coefficients and contractions not expressible in the current support-scoped rule grammar. Exact support is treated as known; sampling-zero discrimination and asymptotics across changing supports remain separate problems.

\section{Honest scope}

The support saturation, district-signature recovery, rational success certificate, rational interior failure certificate, tensor iff characterization, destructive controls, and fixed-face multinomial inference concern the explicitly defined canonical quasi-Markovian class \(\Theta_S\). For every ambient finite-SCM class \(\mathcal M_S\) containing \(\Theta_S\) and preserving the observational and intervention semantics on that subclass, the negative certificate transfers upward, so identification over \(\mathcal M_S\) implies tensor membership. The converse is asserted only for \(\Theta_S\), not for Duarte's broader framework with overlapping latent-parent neighborhoods. The do-search normalization remains open. The diagonal bow and five-cell instrumental-variable controls test the structural-zero boundary inside the canonical class; their negative witnesses, but not their positive conclusions, inherit the ambient-supermodel transfer. Support recovery, optimized bit complexity, and empirical software integration are outside scope.

\begin{thebibliography}{99}

  \bibitem{ShpitserPearl2006} Shpitser, I., and Pearl, J. (2006). Identification of joint interventional distributions in recursive semi-Markovian causal models. \emph{Proceedings of AAAI}, 1219--1226.

  \bibitem{Shpitser2023IDFailure} Shpitser, I. (2023). When does the ID algorithm fail? \emph{arXiv:2307.03750}.

  \bibitem{HwangChoeKwonLee2024} Hwang, I., Choe, Y., Kwon, Y., and Lee, S. (2024). On positivity condition for causal inference. \emph{Proceedings of ICML}, 20818--20841.

  \bibitem{ChenDarwiche2024Constrained} Chen, Y., and Darwiche, A. (2025). Constrained identifiability of causal effects. \emph{arXiv:2412.02869v2}.

  \bibitem{ChenDarwiche2025Granularity} Chen, Y., and Darwiche, A. (2026). On the granularity of causal effect identifiability. \emph{arXiv:2510.16703v3}.

  \bibitem{ZaffalonAntonucciCabanas2020} Zaffalon, M., Antonucci, A., and Cabañas, R. (2020). Structural causal models are solvable by credal networks. \emph{Proceedings of PGM}, 581--592.

  \bibitem{ZaffalonAntonucciCabanasHuberAzzimonti2024} Zaffalon, M., Antonucci, A., Cabañas, R., Huber, D., and Azzimonti, D. (2024). Efficient computation of counterfactual bounds. \emph{International Journal of Approximate Reasoning}, article 109111.

  \bibitem{DuarteFinkelsteinKnoxMummoloShpitser2023} Duarte, G., Finkelstein, N., Knox, D., Mummolo, J., and Shpitser, I. (2023). An automated approach to causal inference in discrete settings. \emph{Journal of the American Statistical Association}.

  \bibitem{ShridharanIyengar2023} Shridharan, M., and Iyengar, G. (2023). Causal bounds in quasi-Markovian graphs. \emph{Proceedings of ICML}.

  \bibitem{TikkaHyttinenKarvanen2019CSI} Tikka, S., Hyttinen, A., and Karvanen, J. (2019). Identifying causal effects via context-specific independence relations. \emph{Advances in Neural Information Processing Systems}.

  \bibitem{TikkaHyttinenKarvanen2021} Tikka, S., Hyttinen, A., and Karvanen, J. (2021). Causal effect identification from multiple incomplete data sources: a general search-based approach. \emph{Journal of Statistical Software}, 99(5), 1--40.

  \bibitem{BhattacharyaNabiShpitser2022} Bhattacharya, R., Nabi, R., and Shpitser, I. (2022). Semiparametric inference for causal effects in graphical models with hidden variables. \emph{Journal of Machine Learning Research}, 23, 1--76.

  \bibitem{vanDerVaart1998} van der Vaart, A. W. (1998). \emph{Asymptotic Statistics}. Cambridge University Press.

\end{thebibliography}

\end{document}